\documentclass[11pt]{amsart}

\usepackage[T1]{fontenc}
\usepackage{lmodern}
\usepackage{microtype}
\usepackage{mathtools,amssymb,amsthm}
\usepackage[hidelinks]{hyperref}

\hypersetup{
  pdftitle={A Double-Suspension Counterexample to the Composition Converse for Strong CW-Regular Subdivisions}
}

\newtheorem{theorem}{Theorem}
\newtheorem{lemma}[theorem]{Lemma}
\newtheorem{proposition}[theorem]{Proposition}
\newtheorem{corollary}[theorem]{Corollary}
\theoremstyle{definition}
\newtheorem*{question}{Question 8.6 of \cite{Stapledon}}
\theoremstyle{remark}
\newtheorem*{remark}{Remark}

\DeclareMathOperator{\rank}{rank}
\newcommand{\face}{\operatorname{face}}
\newcommand{\lk}{\operatorname{lk}}
\newcommand{\Ohat}{\widehat{0}}
\newcommand{\bd}{\partial}
\newcommand{\join}{*}

\title[A double-suspension counterexample]
{A Double-Suspension Counterexample to the Composition Converse for Strong CW-Regular Subdivisions}

\date{}

\subjclass[2020]{06A07, 57Q05, 05E45, 52B22}
\keywords{lower Eulerian poset, strong formal subdivision, CW-regular subdivision,
CW-poset, Poincar\'e homology sphere, double suspension, counterexample}

\begin{document}

\begin{abstract}
Question~8.6 of Stapledon's \emph{Subdivisions of lower Eulerian posets}
(arXiv:2511.16608v1) asks whether the converse of his Lemma 7.15 survives the removal of
the hypothesis $\rank(\sigma)=0$: if $\sigma\colon X\to Y$ and $\tau\colon Y\to Z$ are
strong formal subdivisions of lower Eulerian posets and both $\sigma$ and $\tau\circ\sigma$
are strong CW-regular subdivisions, must $\tau$ be one? We answer \textbf{no} with one
finite triple, of ranks $\rank(\sigma)=2$ and $\rank(\tau)=4$. Let $P$ be the $16$-vertex
triangulation of the Poincar\'e homology $3$-sphere of Bj\"orner and Lutz, $Y=\face(P)$
with the natural rank shifted by $2$, $Z=B_0$ with $\rho_Z(\Ohat_Z)=6$,
$X=\face\bigl(P\join\bd\Delta^2\bigr)$ with its natural rank, $\sigma$ the join projection
and $\tau$ the unique map. Then $|K_X|$ is the double suspension of $|P|$, homeomorphic to
$S^5$ by Cannon's theorem, so $\tau\circ\sigma$ is strong CW-regular; $\sigma$ is strong
CW-regular by the join-of-balls example of the same paper; and $\tau$ is not, because $|P|$
is not a sphere. Nothing beyond this single triple is claimed.
\end{abstract}

\maketitle

\section{The question}\label{sec:question}

The last of the five questions of Section~8 of Stapledon \cite{Stapledon},
\emph{Generalizations and further questions}, reads as follows.

\begin{question}
Does the converse statement in Lemma 7.15 hold without the assumption that $\sigma$ has
rank $0$ (c.f.\ Lemma 3.14)? That is, let $\sigma\colon X\to Y$ and $\tau\colon Y\to Z$ be
strong formal subdivisions between lower Eulerian posets $X,Y,Z$ with rank functions
$\rho_X,\rho_Y,\rho_Z$ respectively. Assume that $\sigma$ and $\tau\circ\sigma$ are strong
CW-regular subdivisions. Is $\tau$ necessarily a strong CW-regular subdivision?
\end{question}

The question is universally quantified over triples, so a single triple decides it. Four
conventions of \cite{Stapledon} are load-bearing.

\begin{enumerate}
\item[(D1)] Definition 7.2: ``We say that $K$ is a regular CW-sphere or regular CW-ball if
  $|K|$ is \emph{homeomorphic} to $S^{\dim K}$ or $D^{\dim K}$ respectively.''
  Homeomorphism only --- not piecewise linear, not shellable.
\item[(D2)] Definition 7.11: an order-preserving, rank-increasing, strongly surjective
  $\sigma\colon X\to Y$ of CW-posets is a \emph{strong CW-regular subdivision} if for every
  $y\in Y$,
  \[
    \dim K_{X,\le y}=\rho_Y(y)-\rho_X(\Ohat_X)-1,
  \]
  and $K_{X,\le y}$ is a regular CW-sphere when $y=\Ohat_Y$, and a regular CW-ball with
  boundary $K_{X,<y}$ otherwise; here $X_{\le y}=\sigma^{-1}\bigl([\Ohat_Y,y]\bigr)$.
\item[(D3)] With the conventions of \cite[\S7.1]{Stapledon}, $B_0$ is a CW-poset, equal to
  the face poset of the empty complex, and $\face(K)$ carries the natural rank
  $\rho(e)=\dim e+1$; the rank of a strong CW-regular subdivision is
  $\rank(\sigma)=\rho_Y(\Ohat_Y)-\rho_X(\Ohat_X)$ (\cite[\S7.2]{Stapledon}).
\item[(D4)] In Lemma 7.15 the forward implication holds unconditionally, and the converse
  is proved \emph{only} under the extra hypothesis $\rank(\sigma)=0$. It is stated for
  CW-posets, which is stronger than the ``lower Eulerian'' of Question~8.6, so a CW-poset
  witness refutes both readings.
\end{enumerate}

\section{The witness}\label{sec:witness}

Let $P$ be the $16$-vertex triangulation of the Poincar\'e homology $3$-sphere
$\Sigma(2,3,5)$ of Bj\"orner and Lutz \cite{BL}, listed by its $90$ tetrahedra in
Table~\ref{tab:P}. It has $f(P)=(16,106,180,90)$ and $\dim P=3$.

\begin{table}[htbp]
\small
\setlength{\tabcolsep}{9pt}
\begin{tabular}{lllll}
\hline
1 2 4 9 & 1 2 4 15 & 1 2 6 14 & 1 2 6 15 & 1 2 9 14 \\
1 3 4 12 & 1 3 4 15 & 1 3 7 10 & 1 3 7 12 & 1 3 10 15 \\
1 4 9 12 & 1 5 6 13 & 1 5 6 14 & 1 5 8 11 & 1 5 8 13 \\
1 5 11 14 & 1 6 13 15 & 1 7 8 10 & 1 7 8 11 & 1 7 11 12 \\
1 8 10 13 & 1 9 11 12 & 1 9 11 14 & 1 10 13 15 & 2 3 5 10 \\
2 3 5 11 & 2 3 7 10 & 2 3 7 13 & 2 3 11 13 & 2 4 9 13 \\
2 4 11 13 & 2 4 11 15 & 2 5 8 11 & 2 5 8 12 & 2 5 10 12 \\
2 6 10 12 & 2 6 10 14 & 2 6 12 15 & 2 7 9 13 & 2 7 9 14 \\
2 7 10 14 & 2 8 11 15 & 2 8 12 15 & 3 4 5 14 & 3 4 5 15 \\
3 4 12 14 & 3 5 10 15 & 3 5 11 14 & 3 7 12 13 & 3 11 13 14 \\
3 12 13 14 & 4 5 6 7 & 4 5 6 14 & 4 5 7 15 & 4 6 7 11 \\
4 6 10 11 & 4 6 10 14 & 4 7 11 15 & 4 8 9 12 & 4 8 9 13 \\
4 8 10 13 & 4 8 10 14 & 4 8 12 14 & 4 10 11 13 & 5 6 7 13 \\
5 7 9 13 & 5 7 9 15 & 5 8 9 12 & 5 8 9 13 & 5 9 10 12 \\
5 9 10 15 & 6 7 11 12 & 6 7 12 13 & 6 10 11 12 & 6 12 13 15 \\
7 8 10 14 & 7 8 11 15 & 7 8 14 15 & 7 9 14 15 & 8 12 14 15 \\
9 10 11 12 & 9 10 11 16 & 9 10 15 16 & 9 11 14 16 & 9 14 15 16 \\
10 11 13 16 & 10 13 15 16 & 11 13 14 16 & 12 13 14 15 & 13 14 15 16 \\
\hline
\end{tabular}
\medskip
\caption{The $90$ tetrahedra of $P$, the $16$-vertex triangulation of the Poincar\'e
homology $3$-sphere \cite{BL}. The properties of $P$ that the proof below consumes ---
its $f$-vector, that it is a pure closed $3$-manifold, that $H_*(P)\cong H_*(S^3)$ over
$\mathbb Z$, and that $|P|$ is not simply connected --- are checked as described in
\S\ref{sec:verify}.}
\label{tab:P}
\end{table}

Put
\begin{align*}
 Y&=\face(P), & \rho_Y&=(\dim{}+1)+2,\ \text{ so }\ \rho_Y(\Ohat_Y)=2,\\
 Z&=B_0=\face(\emptyset), & \rho_Z(\Ohat_Z)&=6,\\
 K_X&=P\join\bd\Delta^2, & \dim K_X&=3+1+1=5,\\
 X&=\face(K_X)=Y\times\bd B_3, & \rho_X&=\dim{}+1,\ \text{ so }\ \rho_X(\Ohat_X)=0,
\end{align*}
with cardinalities $|Y|=393$, $|Z|=1$ and $|X|=7\cdot393=2751$, with
$\sigma\colon X\to Y$ the join projection $(F,G)\mapsto F$ and $\tau\colon Y\to Z$
the unique map. The shift of $\rho_Y$ by $+2$ is licensed by
\cite[Example 3.4]{Stapledon}. Then
\[
 \rank(\sigma)=2-0=2,\quad \rank(\tau)=6-2=4,\quad \rank(\tau\circ\sigma)=6-0=6=2+4 ,
\]
and the two routes to $\rho_Z(\Ohat_Z)$ agree: $0+6=6$ and $2+4=6$. Since
$\rank(\sigma)=2>0$ the triple lies outside the half of Lemma 7.15 that is proved, so it
answers Question~8.6 and contradicts nothing published.

\begin{theorem}\label{thm:main}
For the triple above, $\sigma$ and $\tau$ are strong formal subdivisions, $\sigma$ and
$\tau\circ\sigma$ are strong CW-regular subdivisions, and $\tau$ is not. Hence the
answer to Question~8.6 is \textbf{no}.
\end{theorem}

\begin{proof}
\emph{$\sigma$ is order-preserving, rank-increasing and strongly surjective.} It is a
projection of a product order, hence order-preserving; $\rho_X\bigl((y,z)\bigr)
=(\dim e_y+1)+(\dim e_z+1)\le\rho_Y^{\mathrm{nat}}(y)+2=\rho_Y(\sigma(y,z))$ because
$\rho_{\bd B_3}(z)\in\{0,1,2\}$; and given $(F,G)$ and $y\ge F$ one takes $(y,G')$ with
$G'$ an edge of $\bd\Delta^2$ containing $G$, whose rank is $|y|+2=\rho_Y(y)$, which is
strong surjectivity (\cite[Definition 3.1]{Stapledon}). For
$\tau$, strong surjectivity says every face of $P$ lies in a tetrahedron, and $P$ is
pure, being a triangulated closed $3$-manifold.

\emph{(i) Both maps are strong formal subdivisions,} verified directly from
\cite[Definition 3.5]{Stapledon} rather than deduced from CW-regularity, which would be
circular because (D2) presupposes it. For $\sigma$, the summand
$(-1)^{\rho_Y(y)-\rho_X(x')}$ over $x'=(y,G')\ge(F,G)$ with $\sigma(x')=y$ reduces to
$(-1)^{|G'|}$, and the sum over $G'\supseteq G$ inside $\face(\bd\Delta^2)$ is
$1-3+3=1$ for $G=\emptyset$, $-1+2=1$ for $G$ a vertex and $+1=1$ for $G$ an edge. For
$\tau$ the summand reduces to $(-1)^{|F'|}$, giving $1-16+106-180+90=1$ for
$F=\emptyset$ and, for $|F|=k\ge1$ with $\lk(F)\cong S^{3-k}$,
$(-1)^k(-1)^{4-k}=+1$. All five sums equal $1$. The $\tau$ half is exactly
\cite[Example 7.17]{Stapledon} at $r=4$ with $K=P$.

\emph{(ii) $\sigma$ is a strong CW-regular subdivision,} the only input being the ball
structure of \cite[Example 7.10]{Stapledon}, which we quote. Since
$X_{\le y}=[\Ohat_Y,y]\times\bd B_3$, we have
$K_{X,\le y}=\overline{e_y}\join\bd\Delta^2$, of dimension $\dim e_y+2$, and the
requirement of (D2) is $\rho_Y(y)-\rho_X(\Ohat_X)-1=(\dim e_y+3)-0-1=\dim e_y+2$: equal
for every $y$, including $y=\Ohat_Y$, where $-1+2=1=\dim\bd\Delta^2$. At $y=\Ohat_Y$ one
has $K_{X,\le\Ohat_Y}=\bd\Delta^2$, a triangle boundary, whose space is $S^1$: a regular
CW-sphere. For $y\ne\Ohat_Y$, $\overline{e_y}$ is a simplex $\Delta^d$ and
\cite[Example 7.10]{Stapledon} makes
$\Delta^d\join\bd\Delta^2$ a ball of dimension $d+2$ whose boundary is
$\bd\Delta^d\join\bd\Delta^2$, which is precisely $K_{X,<y}$ because
$X_{<y}=[\Ohat_Y,y)\times\bd B_3$.

\emph{(iii) $\tau\circ\sigma$ is a strong CW-regular subdivision, of rank $6$.} The only
element of $Z$ is $\Ohat_Z$, so $K_{X,\le\Ohat_Z}=K_X$, and the dimension requirement is
$6-0-1=5=\dim K_X$. For the sphere clause, the join of the underlying spaces gives
\[
 |K_X|=|P|\join\bigl|\bd\Delta^2\bigr|\cong|P|\join S^1\cong|P|\join S^0\join S^0
 \cong\Sigma^2|P| ,
\]
using the standard identifications $|\bd\Delta^2|\cong S^1\cong S^0\join S^0$, the
associativity of the join of compact spaces, and $\Sigma W\cong S^0\join W$; so $|K_X|$
is the double suspension of $|P|$. Now $P$ is a closed $3$-manifold with
$H_*(P)\cong H_*(S^3)$ over $\mathbb Z$ --- both checked as described in
\S\ref{sec:verify} --- which is the homology-$3$-sphere hypothesis under which we invoke
Cannon's double suspension theorem \cite{Cannon} (see also \cite{Edwards}); it gives
$\Sigma^2|P|\cong S^5$. By (D1) only a homeomorphism is asked for, so $K_X$ is a regular
CW-sphere.

\emph{(iv) $\tau$ is not a strong CW-regular subdivision.} Again the only element of $Z$
is $\Ohat_Z$ and $K_{Y,\le\Ohat_Z}=P$. The dimension clause \emph{passes}:
$\rho_Z(\Ohat_Z)-\rho_Y(\Ohat_Y)-1=6-2-1=3=\dim P$, so the failure is not a bookkeeping
artefact. The sphere clause fails, because $\pi_1(|P|)\ne1$ by
Proposition~\ref{prop:pi1} below and a space homeomorphic to $S^3$ is simply connected.
Equivalently, by \cite[Example 7.16]{Stapledon} a strong CW-regular subdivision onto
$B_0$ forces $K_Y$ to be a regular CW-sphere.

All hypotheses of Question~8.6 hold and its conclusion fails.
\end{proof}

\section{That $|P|$ is not a sphere, exhibited}\label{sec:pi1}

Clause (iv) is the only place where the failure lives, and it rests on a single negative
fact about $P$, which is exhibitable on the page.

Let $A_5$ be the alternating group on $\{1,\dots,5\}$, written multiplicatively with
left-to-right composition, so that $(gh)(i)=h(g(i))$ and the product of the labels along
an edge path is read left to right. Table~\ref{tab:labels} assigns to each ordered edge
$uv$ of $P$ with $u<v$ an element $g_{uv}\in A_5$, written as the one-line word
$g(1)g(2)g(3)g(4)g(5)$; every edge of $P$ absent from the table carries the identity, and
$g_{vu}:=g_{uv}^{-1}$.

\begin{table}[htbp]
\small
\setlength{\tabcolsep}{9pt}
\begin{tabular}{llll}
\hline
2\,3: 24153 & 2\,5: 35214 & 2\,7: 24153 & 2\,8: 35214 \\
2\,10: 24153 & 2\,11: 35214 & 2\,12: 43521 & 2\,13: 41532 \\
3\,5: 23451 & 3\,11: 23451 & 3\,13: 54213 & 3\,14: 23451 \\
4\,5: 23451 & 4\,6: 23451 & 4\,7: 35214 & 4\,8: 41532 \\
4\,10: 41532 & 4\,11: 35214 & 4\,13: 41532 & 4\,14: 23451 \\
5\,7: 43521 & 5\,9: 25413 & 5\,10: 51234 & 5\,12: 25413 \\
5\,15: 51234 & 6\,7: 43521 & 6\,10: 24153 & 6\,11: 43521 \\
6\,12: 43521 & 7\,9: 31524 & 7\,13: 54213 & 7\,14: 31524 \\
7\,15: 43152 & 8\,9: 25413 & 8\,12: 25413 & 8\,14: 31524 \\
8\,15: 43152 & 9\,10: 35421 & 9\,13: 41532 & 9\,15: 35421 \\
10\,11: 54132 & 10\,12: 54132 & 10\,14: 31524 & 10\,16: 54132 \\
11\,13: 35421 & 11\,15: 43152 & 12\,13: 54213 & 12\,14: 23451 \\
12\,15: 54213 & 13\,14: 54132 & 13\,16: 54132 & 14\,15: 35421 \\
15\,16: 54132 & & & \\
\hline
\end{tabular}
\medskip
\caption{An $A_5$-valued flat connection on the $2$-skeleton of $P$: the $53$ edges of
$P$ whose label is not the identity, each label a $5$-cycle written as a one-line word.
The remaining $53$ edges carry the identity, and they contain a spanning tree of the
$1$-skeleton: for instance the $14$ edges of the star of the vertex~$1$ together with
$14\,16$, which is why the image of the holonomy homomorphism is generated by the labels
themselves.}
\label{tab:labels}
\end{table}

\begin{proposition}\label{prop:pi1}
The labelling of Table~\ref{tab:labels} satisfies
$g_{uv}\,g_{vw}=g_{uw}$ at every one of the $180$ triangles $\{u<v<w\}$ of $P$. Hence
holonomy is a well-defined homomorphism $h\colon\pi_1(|P|)\to A_5$; its image is all of
$A_5$, and the closed edge path $1\to2\to3\to1$ has holonomy
$g_{12}\,g_{23}\,g_{13}^{-1}=(1\,2\,4\,5\,3)\ne1$. In particular $\pi_1(|P|)\ne1$ and
$|P|$ is not homeomorphic to $S^3$, even though $H_*(P)\cong H_*(S^3)$ with no torsion.
\end{proposition}

\begin{proof}
$\pi_1$ of a simplicial complex is the edge-path group of its $2$-skeleton, generated by
the edges outside a spanning tree and presented by one relation
$g_{uv}g_{vw}=g_{uw}$ per triangle. A labelling satisfying every triangle relation is
therefore exactly a homomorphism to $A_5$ once the tree labels are the identity, which
they are; the holonomy of the loop obtained by closing an edge $uv$ through the tree is
$g_{uv}$, so the image of $h$ is the subgroup generated by the labels. The three edges
$12,23,13$ all lie in $P$ while $\{1,2,3\}$ is \emph{not} a face of $P$ --- which is why
its holonomy is not forced to be trivial --- and $g_{12}=g_{13}=1$ while
$g_{23}=(1\,2\,4\,5\,3)$. The $180$ relations, the generated image and the homology of
$P$ are checked as described in \S\ref{sec:verify}; every one of them is a finite
computation in exact arithmetic on Table~\ref{tab:P} and Table~\ref{tab:labels}.
\end{proof}

\section{Verification}\label{sec:verify}

The two objects the argument consumes are printed above in full: the $90$ tetrahedra of
Table~\ref{tab:P} and the $53$ labels of Table~\ref{tab:labels}. The accompanying program
\texttt{verify.py} (Python 3.9+, standard library only, exact integer arithmetic, no
external data) transcribes exactly those two tables and re-derives the finite quantities
claimed here: $f(P)=(16,106,180,90)$; that every triangle of $P$ lies in exactly two
tetrahedra, every vertex link is a combinatorial $2$-sphere and every edge link a single
cycle, so $P$ is a closed $3$-manifold; that $H_*(P)=(\mathbb Z,0,0,\mathbb Z)$ with no
torsion, by Smith normal form over $\mathbb Z$; the $180$ triangle relations of
Proposition~\ref{prop:pi1}, that the labels generate a group of order $60$, and the
nontrivial holonomy of $1\to2\to3\to1$; $f(K_X)=(19,157,546,948,810,270)$ and
$|X|=2751=393\cdot7$; the rank arithmetic $2+4=6$ from both sides; that $\sigma$ is
order-preserving, rank-increasing and strongly surjective; all $23{,}359$ alternating sums
of \cite[Definition 3.5]{Stapledon} for $\sigma$ and all $393$ for $\tau$, each equal to
$1$; and the dimension clause of (D2) at every one of the $393$ elements of $Y$ together
with the boundary bookkeeping at the $392$ elements $y\ne\Ohat_Y$. The program reports
\texttt{VERDICT: ALL 52 CHECKS PASS}.

Three inputs to the proof are quoted rather than reproved, and the program does not touch
them: Cannon's theorem, and with it every homeomorphism type, none of which is computed
anywhere; the ball structure of \cite[Example 7.10]{Stapledon}, of which only the
combinatorial consequences are checked; and the classical identification of $\pi_1(|P|)$,
of which only the surjection onto $A_5$ is proved and only that is needed. This paper does
not address the case $\rank(\sigma)=1$ of Question~8.6; it claims no minimality, and it
makes no claim about other complexes in place of $P$, about other ranks, or about the
literature, which the program does not search.

\begin{thebibliography}{9}

\bibitem{BL}
A.~Bj\"orner and F.~H. Lutz,
\emph{Simplicial manifolds, bistellar flips and a $16$-vertex triangulation of the
Poincar\'e homology $3$-sphere},
Experiment. Math. \textbf{9} (2000), no.~2, 275--289.
\href{https://doi.org/10.1080/10586458.2000.10504652}{doi:10.1080/10586458.2000.10504652};
MR1780212; Zbl 1101.57306.

\bibitem{Cannon}
J.~W. Cannon,
\emph{Shrinking cell-like decompositions of manifolds. Codimension three},
Ann. of Math. (2) \textbf{110} (1979), no.~1, 83--112.
\href{https://doi.org/10.2307/1971245}{doi:10.2307/1971245}; Zbl 0424.57007.

\bibitem{Edwards}
R.~D. Edwards,
\emph{Suspensions of homology spheres},
\href{https://arxiv.org/abs/math/0610573}{arXiv:math/0610573}, 2006 (circulated 1975).

\bibitem{Stapledon}
A.~Stapledon,
\emph{Subdivisions of lower Eulerian posets},
\href{https://arxiv.org/abs/2511.16608v1}{arXiv:2511.16608v1}, 20 November 2025.

\end{thebibliography}

\end{document}
